\subsection{The Gram matrix when a weight varies}
\hypertarget{source-tied-gram}{}

To apply Lemma~\ref{lem:smith-average}, we need a matrix polynomial with a controlled determinant. A decomposition by the degree of the final vector produces it from the Gram form.

Let \(f\) be a polynomial in the \(b\) vector or covector variables of a fixed initial segment, each having \(r\) coordinates. Their degrees are \(o_i\), the
total degree is \(S\), and the last letter \(w\) is a vector of
degree \(o\). The others may be vectors or covectors. The Fock
pairing is bilinear, with orthogonal monomials and squared
monomial norms \(\prod_\alpha m_\alpha!\). Choose a coordinate
\(\rho\) and add a covector \(\eta\) of variable degree \(n\ge S\).
Define the substitution \(T_\eta\) on each vector \(v\) and covector \(u\). For every \(i\ne\rho\) its coordinate rules are
\[
 \begin{aligned} (T_\eta v)_\rho&=v\cdot\eta, & (T_\eta v)_i&=\eta_\rho v_i,\\ (T_\eta u)_\rho&=\eta_\rho u_\rho, & (T_\eta u)_i&=\eta_\rho u_i-u_\rho\eta_i. \end{aligned}
\]
Set \(f_n=\eta_\rho^{n-S}T_\eta f\). Evaluating at the \(\rho\)-th coordinate vector recovers
\(f\). In particular the substitution is an integral injection.

Let \(U\) multiply by \(w\cdot\eta\), let
\(D=\sum_i\partial_{w_i}\partial_{\eta_i}\), and let \(z\)
be the nonnegative integer decoration of this adjacent pair.
Write \(B=o+n+r+z-1\) and define
\[
 \Phi_z=\sum_{j=0}^o\frac{(-1)^j U^jD^j}
                           {(j!)^2\binom Bj},\qquad
 \Gamma_n(f,g)=\binom Bo\,
       \frac{\langle\Phi_z f_n,g_n\rangle}
            {(n-S)!\,S!\prod_i o_i!}.
\]
Multiplication \(U\) and contraction \(D\) are adjoint for the Fock pairing. The scalar
in this formula is used before the bound on small Smith factors.

\begin{lemma}[Polynomial Gram pairing at a tied pair]
\label{lem:tied-gram-polynomial}
For fixed inputs and fixed decoration \(z\), \(\Gamma_n\) is polynomial
in \(n\) of degree at most \(S+o\). For integral inputs its entries
after multiplication by \(p^u\) lie in the integral binomial basis, with
\[
 u=\left\lceil ((b+3)r+1)
                    (1+\log_p(\max(1,3S)))\right\rceil .
\]
If a basis \((f_i)\) is adapted to the filtration by order in \(w_\rho\),
write \(e_i\) for the order of \(f_i\), \(l\) for its size and
\(d_0=\sum_i(S+o-e_i)\). Suppose each normalized graded determinant
has valuation at most \(v_*\) times its rank. For \(S,z\le2a\) the
determinant of \((p^u\Gamma_n(f_i,f_j))\) has degree \(d_0\) and leading
coefficient \(C\) with
\[
 v(C)+d_0/(p-1)\le
     l\big(u+v_*+3(1+\log_p(8a+r))\big).
\]
\end{lemma}

\begin{proof}
\proofstep{Degree in the varied weight.}
Each entry is polynomial in \(n\) of degree at most \(S+o\).
In the \(j\)-th term move \(U^j\) by adjunction.
The ratio \(\binom Bo/\binom Bj=\binom{B-j}{o-j}/\binom oj\)
is polynomial of degree \(o-j\). To see the remaining polynomial,
put \(m=n-S\). Before applying \(D^j\), the two occurrences of
\(\eta_\rho\) in a monomial pairing have powers \(m+t,m+t'\),
where \(0\le t,t'\le S\). If the differentiation orders there
are \(a,a'\), the residual powers must agree. Their contribution,
including the Fock factorial and division by \(m!\), is
\[
   \frac{(m+t)!}{m!}(m+t')_{\underline{a'}},
       \qquad t-a=t'-a' .
\]
All remaining factors are independent of \(m\). This polynomial
has degree \(\le S+j\), and its falling factorial also gives the
correct zero when the differentiation order is too large.

\proofstep{Integrality in the binomial basis.}
For integral inputs we can clear the binomial-basis denominators
with logarithmic loss. Evaluate at the consecutive integers
\(S\le n\le S+S+o\le3S\). In the expansion of \(D^j/j!\)
let \(h_i\) contractions take place at coordinate \(i\).
On a monomial with exponents \(u_i,v_i\) in \(w,\eta\), the multi-index \(h\) contributes the coefficient
\[
 \prod_i\frac{u_i^{\underline{h_i}}v_i^{\underline{h_i}}}{h_i!} =\prod_i h_i!\binom{u_i}{h_i}\binom{v_i}{h_i},\qquad \sum_i h_i=j.
\]
Use this on both sides, with multiplicities \(h_i,h'_i\).

If the resulting monomial has exponent vector \(d\), the
numerator contains \(\prod_\alpha d_\alpha!\prod_i h_i!h'_i!\).
The sum of the factorial arguments is exactly \(S+n\), the
same as for \((n-S)! S!\prod_i o_i!\). Legendre's formula
gives
\[
  0\le t/(p-1)-v(t!)\le 1+\log_p(\max(1,t)).
\]
There are at most \((b+3)r\) numerator factorials. The additional
denominator \(\binom oj\) costs at most \(1+\log_p(\max(1,o))\)
by counting carries in base \(p\). The other factor
\(\binom{B-j}{o-j}\) is integral for every \(z\). Thus the
values become integral after multiplication by \(p^u\) for
\[
 u=\left\lceil\big((b+3)r+1\big)
                \big(1+\log_p(\max(1,3S))\big)\right\rceil .
\]
Taking forward differences at \(S\) recovers the integral
coefficients in the basis \(\binom{n-S}{j}\). Translation
through an integer preserves the integral binomial basis,
so the hypothesis on the entries in \eqref{eq:smith-average}
holds with \(\delta=2S\).

\proofstep{The leading harmonic pairing.}
We compute the leading coefficient to keep the dependence on
\(S\) out of the final valuation allowance. Write \(e(f)\) for
the least exponent of \(w_\rho\) in a nonzero \(f\), and \(I_e f\)
for its terms of exponent exactly \(e\). For
\(e=e(f),e'=e(g)\) one has
\begin{align}
 \deg\Gamma_n(f,g)&\le S+o-\max(e,e'),\label{eq:tie-degree}\\
 {}[n^{S+o-e}]\Gamma_n(f,g)
   &=\frac{e!\binom{z+e}{e}}{o!\,S!\prod_i o_i!}
                   \langle I_e f,I_e g\rangle
                      \qquad(e=e').\label{eq:tie-top}
\end{align}
It suffices to prove these rational identities on monomials,
so we can use the positive real Fock metric to compute the
leading coefficient there. On bidegree \((o,n)\) the orthogonal
decomposition by powers of \(U\) has summands \(U^j\ker D\)
with harmonic factors of degrees \((o-j,n-j)\). Indeed
\((\ker D)^\perp=\operatorname{im}U\), and iteration gives the
summands. The commutator \([D,U]\) is the sum of the two Euler
operators plus \(r\), which proves orthogonality by moving powers
of \(U\) across the pairing.

It gives the eigenvalue of
\(\Phi_z\) on the \(j\)-th summand as
\[
 \sum_{l=0}^j (-1)^l\binom jl
       \binom{o+n+r-j-1}{l}/\binom Bl
      =\binom{z+j}{j}/\binom Bj .
\]
The equality is inclusion--exclusion for a sample of \(j\)
objects from \(B\) objects, of which \(o+n+r-j-1\) are forbidden.
Thus the eigenvalue is \(O(n^{-j})\) for fixed prefix and \(z\).
The image \(f_n\) belongs to \(U^e\) times the smaller bidegree,
since the substitution sends \(w_\rho\) to \(w\cdot\eta\).
Also
\[
 \|f_n\|^2/(n-S)!=O(n^S)
\]
by expanding its monomials.
Only levels \(j\ge\max(e,e')\) can pair, proving
\eqref{eq:tie-degree}. The real limits here extract coefficients
at fixed prefix. The uniform estimates used as the prefix grows
are the ensuing factorial valuations.

For \(e=e'\), write \(f_n=U^e h_n\). The polynomial \(h_n\)
is \(\eta_\rho^{n-S}\) times the substitution on
\(h=f/w_\rho^e\), and has bidegree \((o-e,n-e)\).
Use \(g_n=U^e h'_n\) in the same way. After division by \((n-S)!\), a level \(j>e\) contributes at most
\[
 O(n^{S})\,O(n^{-j})=O(n^{S-j})=O(n^{S-e-1}).
\]
Thus only level \(e\) can survive at order \(n^{S-e}\).

Let \(P\) be the harmonic projector on this smaller bidegree.
Its formula is the sum defining \(\Phi_z\), with degrees
\(o-e,n-e\) and denominator parameter \(o+n-2e+r-2\).
The same eigenvalue identity gives one on level zero and zero
on the higher levels (use the binomial identity at \(z=-1\)).
Only large \(n\) are needed in this limit. On harmonic vectors
\[
 D^eU^e=e!\,(o+n-2e+r)^{\overline e}.
\]
Multiplying by the level-\(e\) eigenvalue therefore tends to
\((e!)^2\binom{z+e}{e}\).

The remaining harmonic pairing has limit
\[
  \lim_{n\to\infty}\frac{\langle P h_n,h'_n\rangle}
                            {(n-S)!\,n^{S-e}}
       =\langle h|_{w_\rho=0},h'|_{w_\rho=0}\rangle .
\]
In the \(l\)-th term of \(P\), move the \(l\) multiplications
across as derivatives. The term before the projector coefficient,
divided by \((n-S)! n^{S-e+l}\), tends to
\(\langle \partial_{w_\rho}^l h,\partial_{w_\rho}^l h'\rangle\).
To check the powers, let \(d,d'\) count occurrences of
\(\eta_i\), \(i\ne\rho\), and \(a,a'\) the respective numbers
of derivatives falling there. A contributing pair has
\(d-a=d'-a'\). After factorial division its degree is at most
\[
 S-e+l-\underbrace{(d-a)}_{\ge0}-a-a'\le S-e+l.
\]
Equality forces \(a=a'=0\) and \(d-a=d'-a'=0\), hence \(d=d'=0\).

The coefficient of the pure
\(\eta_\rho\)-power in the substituted \(h\) is \(h\) itself.
This proves the limit term by term. The leading projector
coefficient is \((-1)^l/(l! n^l)\). On the terms of \(w_\rho\)-degree \(t\) in each argument,
moving the derivatives across the Fock form and the factor
\(1/l!\) give the scalar
\[
       \sum_{l=0}^t(-1)^l\binom tl
           =\begin{cases}1,&t=0,\\0,&t>0.\end{cases}
\]
Distinct degrees are orthogonal. Thus the limiting operator
keeps only the restrictions \(w_\rho=0\). Restoring the factor
\(w_\rho^e\) in the initial terms multiplies this pairing by
\(e!\). Finally \(\binom Bo/n^o\to1/o!\),
which proves \eqref{eq:tie-top}.

\proofstep{The leading determinant in a filtration basis.}
Choose a basis \(f_i\) adapted to the integral filtration by
vanishing order, so that the \(I_e f_i\) of order \(e\) form
a basis for the image at graded order \(e\). Such a basis
exists over a DVR by successive lifting from the free graded
images of the coefficient map. In the determinant of
\((p^u\Gamma_n(f_i,f_j))\) the degree is at most
\(d_0=\sum_i(S+o-e_i)\). For a permutation \(\pi\) the excess in vanishing orders is
\[
 \sum_i\max(e_i,e_{\pi(i)})-\sum_i e_i =\tfrac12\sum_i|e_i-e_{\pi(i)}|.
\]
Thus its determinant term can attain degree \(d_0\) only if \(e_i=e_{\pi(i)}\) throughout, by
\eqref{eq:tie-degree} and \(\sum_i e_i=\sum_i e_{\pi(i)}\).

Consequently its top coefficient is the product of the
determinants of the graded Fock forms normalized by \(\prod o_i!\),
times
\[
                  \prod_i\frac{p^u e_i!\binom{z+e_i}{e_i}}
                                         {o!\,S!}.
\]
If each normalized graded determinant has valuation at most
\(v_*\) times its rank, the top coefficient is nonzero. For
\(S,z\le2a\), Legendre's formula and the carry bound give
\[
  v(C)+d_0/(p-1)\le l\big(u+v_*+3(1+\log_p(8a+r))\big),
\]
where \(l\) is the size of the matrix. Indeed the three
linear factorial terms cancel \((S+o-e_i)/(p-1)\) in every
row, and the digit error from \(e_i!\) has favorable sign.
This supplies the polynomial and the precise leading-coefficient
input to \eqref{eq:smith-average} from a bound on the graded
form.
\end{proof}
